\documentclass{beamer}


\usepackage[T2A]{fontenc}
\usepackage[utf8]{inputenc}
\usepackage[english,russian]{babel}
\pdfmapfile{+sansmathaccent.map}


\mode<presentation>
{
	\usetheme{Warsaw} % or try Darmstadt, Madrid, Warsaw, Rochester, CambridgeUS, ...
	\usecolortheme{seahorse} % or try seahorse, beaver, crane, wolverine, ...
	\usefonttheme{serif}  % or try serif, structurebold, ...
	\setbeamertemplate{navigation symbols}{}
	\setbeamertemplate{caption}[numbered]
} 


%%%%%%%%%%%%%%%%%%%%%%%%%%%%
% itemize settings


%%%%%%%%%%%%%%%%%%%%%%%%%%%%
% itemize settings

\definecolor{myhotpink}{RGB}{255, 80, 200}
\definecolor{mywarmpink}{RGB}{255, 60, 160}
\definecolor{mylightpink}{RGB}{255, 80, 200}
\definecolor{mypink}{RGB}{255, 30, 80}
\definecolor{mydarkpink}{RGB}{155, 25, 60}

\definecolor{mypaleblue}{RGB}{240, 240, 255}
\definecolor{mylightblue}{RGB}{120, 150, 255}
\definecolor{myblue}{RGB}{90, 90, 255}
\definecolor{mygblue}{RGB}{70, 110, 240}
\definecolor{mycourseblue}{RGB}{120, 140, 255}
\definecolor{mycourselightblue}{RGB}{245, 250, 255}
\definecolor{mydarkblue}{RGB}{0, 0, 180}
\definecolor{myblackblue}{RGB}{40, 40, 120}

\definecolor{mygreen}{RGB}{0, 200, 0}
\definecolor{mygreen2}{RGB}{245, 255, 230}
\definecolor{mydarkgreen}{RGB}{0, 120, 0}
\definecolor{mydarkgreen}{RGB}{0, 120, 0}


\definecolor{mygray}{gray}{0.8}
\definecolor{mydarkgray}{RGB}{80, 80, 160}
\definecolor{mylightgray}{RGB}{160, 160, 160}

\definecolor{mydarkred}{RGB}{160, 30, 30}
\definecolor{mylightred}{RGB}{255, 150, 150}
\definecolor{myred}{RGB}{200, 110, 110}
\definecolor{myblackred}{RGB}{120, 40, 40}

\definecolor{mypink}{RGB}{255, 30, 80}
\definecolor{myhotpink}{RGB}{255, 80, 200}
\definecolor{mywarmpink}{RGB}{255, 60, 160}
\definecolor{mylightpink}{RGB}{255, 80, 200}
\definecolor{mydarkpink}{RGB}{155, 25, 60}
\definecolor{mywhitepink}{RGB}{255, 240, 240}

\definecolor{mydarkcolor}{RGB}{60, 25, 155}
\definecolor{mylightcolor}{RGB}{130, 180, 250}



\setbeamertemplate{itemize items}[default]

\setbeamertemplate{itemize item}{\color{myblackblue}$\blacksquare$}
\setbeamertemplate{itemize subitem}{\color{mygblue}$\blacktriangleright$}
\setbeamertemplate{itemize subsubitem}{\color{mygray}$\blacksquare$}

\setbeamercolor{palette quaternary}{fg=white,bg=mycourseblue}
\setbeamercolor{titlelike}{parent=palette quaternary}

\setbeamercolor{palette quaternary2}{fg=black,bg=mycourselightblue}
\setbeamercolor{frametitle}{parent=palette quaternary2}

\setbeamerfont{frametitle}{size=\Large,series=\scshape}
\setbeamerfont{framesubtitle}{size=\normalsize,series=\upshape}





%%%%%%%%%%%%%%%%%%%%%%%%%%%%
% block settings

\setbeamercolor{block title}{bg=red!30,fg=black}

\setbeamercolor*{block title example}{bg=mygreen!40!white,fg=black}

\setbeamercolor*{block body example}{fg= black, bg= mygreen2}


%%%%%%%%%%%%%%%%%%%%%%%%%%%%
% URL settings
\hypersetup{
	colorlinks=true,
	linkcolor=blue,
	filecolor=blue,      
	urlcolor=blue,
}

%%%%%%%%%%%%%%%%%%%%%%%%%%

\renewcommand{\familydefault}{\rmdefault}


\newcommand{\mycolor}[1] {
\setbeamercolor{normal text}{fg={#1}} % Recolors standard body text
\setbeamercolor{structure}{fg={#1}}   % Recolors bullets and structure items
\usebeamercolor[fg]{normal text}      % Applies the normal text color immediately
}


\usepackage{amsmath}
\usepackage{mathtools}

\usepackage{cancel}
\usepackage{centernot}

\usepackage{changepage}
\usepackage{subcaption}

\usepackage{qrcode}

\newcommand{\bo}[1] {\mathbf{#1}}
\newcommand{\R}{\mathbb{R}} 
\newcommand{\T}{^\top}     



\newcommand{\mydate}{Fall 2026}

\newcommand{\mygit}{\textcolor{blue}{\href{https://github.com/SergeiSa/Advanced-Control-2026}{github.com/SergeiSa/Advanced-Control-2026}}}

\newcommand{\myqr}{ \textcolor{black}{\qrcode[height=1.5in]{https://github.com/SergeiSa/Advanced-Control-2026}}
}




\newcommand{\bref}[2] {\textcolor{blue}{\href{#1}{#2}}}

%%%%%%%%%%%%%%%%%%%%%%%%%%%%
% code settings

\usepackage{listings}
\usepackage{color}
% \definecolor{mygreen}{rgb}{0,0.6,0}
% \definecolor{mygray}{rgb}{0.5,0.5,0.5}
\definecolor{mymauve}{rgb}{0.58,0,0.82}
\lstset{ 
	backgroundcolor=\color{white},   % choose the background color; you must add \usepackage{color} or \usepackage{xcolor}; should come as last argument
	basicstyle=\footnotesize,        % the size of the fonts that are used for the code
	breakatwhitespace=false,         % sets if automatic breaks should only happen at whitespace
	breaklines=true,                 % sets automatic line breaking
	captionpos=b,                    % sets the caption-position to bottom
	commentstyle=\color{mygreen},    % comment style
	deletekeywords={...},            % if you want to delete keywords from the given language
	escapeinside={\%*}{*)},          % if you want to add LaTeX within your code
	extendedchars=true,              % lets you use non-ASCII characters; for 8-bits encodings only, does not work with UTF-8
	firstnumber=0000,                % start line enumeration with line 0000
	frame=single,	                   % adds a frame around the code
	keepspaces=true,                 % keeps spaces in text, useful for keeping indentation of code (possibly needs columns=flexible)
	keywordstyle=\color{blue},       % keyword style
	language=Octave,                 % the language of the code
	morekeywords={*,...},            % if you want to add more keywords to the set
	numbers=left,                    % where to put the line-numbers; possible values are (none, left, right)
	numbersep=5pt,                   % how far the line-numbers are from the code
	numberstyle=\tiny\color{mygray}, % the style that is used for the line-numbers
	rulecolor=\color{black},         % if not set, the frame-color may be changed on line-breaks within not-black text (e.g. comments (green here))
	showspaces=false,                % show spaces everywhere adding particular underscores; it overrides 'showstringspaces'
	showstringspaces=false,          % underline spaces within strings only
	showtabs=false,                  % show tabs within strings adding particular underscores
	stepnumber=2,                    % the step between two line-numbers. If it's 1, each line will be numbered
	stringstyle=\color{mymauve},     % string literal style
	tabsize=2,	                   % sets default tabsize to 2 spaces
	title=\lstname                   % show the filename of files included with \lstinputlisting; also try caption instead of title
}


%%%%%%%%%%%%%%%%%%%%%%%%%%%%
% URL settings
\hypersetup{
	colorlinks=false,
	linkcolor=blue,
	filecolor=blue,      
	urlcolor=blue,
}

%%%%%%%%%%%%%%%%%%%%%%%%%%

%%%%%%%%%%%%%%%%%%%%%%%%%%%%
% tikz settings

\usepackage{tikz}
\tikzset{every picture/.style={line width=0.75pt}}


\newcommand{\myqrframe}{
	\begin{frame}
		\begin{center}
			\centerline{Слайды доступны на Github:}
			\bigskip
			\centerline{\mygit}
			\bigskip
			\myqr
		\end{center}
	\end{frame}
}


\subtitle{Современные методы управления роботами}
\author{Сергей Савин}
\date{2026}

\title{Планирование пути на графах}
\begin{document}
	\maketitle
	
	
	\begin{frame}{Содержание}
		\begin{itemize}
			\item Планирование пути на графах
			\item Алгоритм Дейкстры: инициализация, исследование и обновление узлов, выбор следующего узла, восстановление пути
			\item Алгоритм A$^\ast$ и допустимые эвристики
			\item Итерация по значениям и уравнение Беллмана
		\end{itemize}
	\end{frame}
	
	
	%------------------------------------------------
	\begin{frame}{Планирование пути на графах}
		\begin{flushleft}
			
			Многие практические задачи планирования пути можно
			формализовать с помощью графов:
			
			\begin{itemize}
				\item Мы описываем среду как набор
				местоположений (узлов).
				\item Мы знаем, между какими местоположениями можно
				перемещаться (рёбра).
				\item Мы знаем стоимость перемещения из одного
				местоположения в другое (стоимость ребра).
			\end{itemize}
			
			Таким образом, планирование пути становится задачей поиска
			последовательности рёбер, соединяющих начальный узел
			с целевым узлом, такой, что сумма стоимостей рёбер
			минимальна.
			
		\end{flushleft}
	\end{frame}
	
	%------------------------------------------------
	\begin{frame}{Пути на графе}
		\begin{flushleft}
			
			Для любого узла $k$ пусть $\mathcal{N}(k)$ обозначает
			множество узлов, достижимых из $k$
			прохождением одного ребра. Они называются
			\textbf{смежными узлами}.
			
			\medskip
			
			Путь можно представить как последовательность узлов,
			соединённых рёбрами. Каждый промежуточный узел имеет
			предшественника и преемника.
			
			
			\medskip
			
			
			Оптимальный путь обладает важным свойством:
			
			\begin{itemize}
				\item Каждый промежуточный узел достигается из
				своего предшественника вдоль оптимального пути к этому
				узлу.
				\item Из каждого промежуточного узла оставшаяся
				часть пути к цели также оптимальна.
			\end{itemize}
			
			
		\end{flushleft}
	\end{frame}
	
	%------------------------------------------------
	\begin{frame}{Алгоритм Дейкстры}
		\begin{flushleft}
			
			\textbf{Алгоритм Дейкстры} находит кратчайшие пути
			от начального узла к другим узлам графа.
			Его также можно использовать для поиска кратчайшего пути к
			конкретному целевому узлу.
			
			
			\medskip
			Ключевая идея — поддерживать наилучшую известную
			стоимость достижения каждого обнаруженного узла из
			начального узла.
			
			\medskip
			
			Мы поддерживаем два множества:
			
			\begin{itemize}
				\item \textbf{Открытое множество:} узлы, которые были
				обнаружены, но ещё не полностью исследованы.
				\item \textbf{Закрытое множество:} узлы, стоимости кратчайших
				путей до которых уже установлены.
			\end{itemize}
			
			Мы предполагаем, что все стоимости рёбер неотрицательны.
			
		\end{flushleft}
	\end{frame}
	
	%------------------------------------------------
	\begin{frame}{Алгоритм Дейкстры: инициализация}
		\begin{flushleft}
			
			Изначально обнаружен только начальный узел $s$:
			
			\[
			\text{Open}=\{s\},
			\qquad
			\text{Closed}=\emptyset.
			\]
			
			Для каждого обнаруженного узла $k$ мы поддерживаем
			предварительную стоимость $g(k)$: стоимость лучшего пути
			к $k$, найденного на данный момент. Начальная стоимость:
			\[
			g(s)=0.
			\]
			
			Все остальные узлы изначально не обнаружены, их
			стоимости считаются бесконечными:
			\[
			g(k)=+\infty,\qquad k\ne s.
			\]
			
			Мы также сохраняем предшественника каждого обнаруженного
			узла для восстановления пути.
			
		\end{flushleft}
	\end{frame}
	
	%------------------------------------------------
	\begin{frame}{Алгоритм Дейкстры: исследование узла}
		\begin{flushleft}
			
			Исследование узла $k$ включает следующие шаги:
			
			\begin{enumerate}
				\item Проверить, является ли $k$ целевым узлом.
				\item Рассмотреть каждый смежный узел $m\in\mathcal{N}(k)$.
				\item Переместить $k$ из открытого множества в
				закрытое.
			\end{enumerate}
			
			
			\medskip
			При рассмотрении смежного узла $m$ мы проверяем,
			улучшает ли путь через $k$ его
			наилучшую известную стоимость. Если $m$ не был обнаружен, его предварительная
			стоимость инициализируется как
			\[
			g(m)=g(k)+c(k,m),
			\]
			где $c(k,m)$ — стоимость ребра из
			$k$ в $m$.
			
			Мы также записываем $k$ как предшественника $m$
			и добавляем $m$ в открытое множество.
			
		\end{flushleft}
	\end{frame}
	
	%------------------------------------------------
	\begin{frame}{Алгоритм Дейкстры: обновление узла}
		\begin{flushleft}
			
			Иногда мы встречаем узел $m$, который уже
			был обнаружен.
			
			Пусть $g(m)$ — его текущая предварительная стоимость. Мы
			сравниваем её со стоимостью достижения $m$ через
			узел $k$, который сейчас исследуется:
			
			\[
			g_{\mathrm{new}}(m)=g(k)+c(k,m).
			\]
			
			Возможны два варианта:
			
			\begin{itemize}
				\item Если $g(m)\leq g_{\mathrm{new}}(m)$,
				сохранить текущую стоимость и предшественника.
				\item Если $g(m)>g_{\mathrm{new}}(m)$,
				обновить стоимость и назначить $k$ новым
				предшественником $m$:
				\[
				g(m)\leftarrow g_{\mathrm{new}}(m).
				\]
			\end{itemize}
			
			Эта операция называется \textbf{релаксацией}
			ребра.
			
		\end{flushleft}
	\end{frame}
	
	%------------------------------------------------
	\begin{frame}{Алгоритм Дейкстры: выбор следующего узла}
		\begin{flushleft}
			
			Ключевой шаг — выбор следующего узла для исследования.
			
			Мы выбираем узел $k$ с наименьшей предварительной
			стоимостью среди всех узлов открытого множества:
			
			\[
			k=\arg\min_{m\in\mathrm{Open}} g(m).
			\]
			
			Затем мы исследуем $k$ и перемещаем его в закрытое множество.
			
			\begin{itemize}
				\item Если $k$ — целевой узел, мы завершаем работу.
				\item В противном случае мы продолжаем выбирать и
				исследовать узлы.
			\end{itemize}
			
			Поскольку все стоимости рёбер неотрицательны,
			стоимость выбранного узла является стоимостью его кратчайшего пути.
			
		\end{flushleft}
	\end{frame}
	
	
	
	%------------------------------------------------
	\begin{frame}{Алгоритм Дейкстры: восстановление пути}
		\begin{flushleft}
			
			После завершения работы алгоритма мы можем восстановить
			кратчайший путь к цели. Начиная с целевого узла, мы многократно находим
			его предшественника и добавляем его к пути.
			
			\medskip
			
			Эквивалентно, для текущего узла $m$ мы можем
			выбрать предшественника $k$, удовлетворяющего условию
			\[
			k=\arg\min_{k:\,m\in\mathcal{N}(k)}
			\bigl(g(k)+c(k,m)\bigr).
			\]
			
			Мы продолжаем, пока не достигнем начального узла $s$.
			Затем мы обращаем последовательность узлов, чтобы получить
			путь от $s$ к цели.
			
		\end{flushleft}
	\end{frame}
	
	%------------------------------------------------
	\begin{frame}{Алгоритм Дейкстры: одна цель или все цели?}
		\begin{flushleft}
			
			Алгоритм Дейкстры использует целевой узел только
			для определения момента завершения.
			
			
			\medskip
			
			Если мы позволим алгоритму работать до тех пор, пока открытое множество
			не станет пустым, он вычислит стоимости кратчайших путей от
			начального узла до каждого достижимого
			узла графа.
			
			\begin{itemize}
				\item Затем мы можем выбрать любой достижимый узел
				в качестве цели и восстановить его кратчайший путь.
				\item Однако это может быть излишней работой,
				если нам нужен путь только к одной известной цели.
			\end{itemize}
			
			Мы можем использовать знание о цели, чтобы направлять
			поиск и исследовать меньше узлов.
			
		\end{flushleft}
	\end{frame}
	
	%------------------------------------------------
	\begin{frame}{Алгоритм A$^\ast$}
		\begin{flushleft}
			
			Алгоритм A$^\ast$ похож на алгоритм Дейкстры,
			но использует информацию о цели
			для направления выбора узлов для исследования.
			
			\medskip
			В то время как Дейкстра выбирает узел с наименьшей
			предварительной стоимостью $g(k)$, A$^\ast$ выбирает узел
			с наименьшей оценочной полной стоимостью:
			
			\[
			f(k)=g(k)+h(k).
			\]
			
			Здесь:
			\begin{itemize}
				\item $g(k)$ — стоимость лучшего пути,
				найденного на данный момент от начала до узла $k$.
				\item $h(k)$ оценивает оставшуюся стоимость
				от узла $k$ до цели.
			\end{itemize}
			
			\textbf{Допустимая эвристика} никогда не переоценивает
			истинную минимальную оставшуюся стоимость.
			
		\end{flushleft}
	\end{frame}
	
	
	
	
	%------------------------------------------------
	\begin{frame}{Итерация по значениям}
		\begin{flushleft}
			
			Алгоритм Дейкстры начинается с начального узла
			и распространяет информацию о стоимости
			достижения других узлов.
			
			
			\medskip
			
			В качестве альтернативы мы можем начать с цели и
			распространять информацию о стоимости достижения
			цели из каждого другого узла.
			
			\bigskip
			Эта величина называется \textbf{стоимостью до цели}
			( cost-to-go ) и обозначается $V(k)$ для каждого узла $k$. Это ключевая идея \textbf{итерации по значениям}.
			
		\end{flushleft}
	\end{frame}
	
	%------------------------------------------------
	\begin{frame}{Итерация по значениям: инициализация}
		\begin{flushleft}
			
			Мы инициализируем каждый узел $k$ оценкой стоимости до цели
			$V(k)$.
			
			\begin{itemize}
				\item Для целевого узла $g$ стоимость до цели
				равна нулю:
				\[
				V(g)=0.
				\]
				
				\item Для всех остальных узлов начальная оценка
				равна бесконечности:
				\[
				V(k)=+\infty,\qquad k\ne g.
				\]
			\end{itemize}
			
			Затем мы итеративно обновляем эти оценки, используя
			стоимости перемещения к смежным узлам.
			
		\end{flushleft}
	\end{frame}
	
	%------------------------------------------------
	\begin{frame}{Итерация по значениям: обновления}
		\begin{flushleft}
			
			Мы обновляем оценки стоимости до цели, используя
			\textbf{уравнение Беллмана}:
			
			\[
			V(k)=\min_{m\in\mathcal{N}(k)}
			\bigl(c(k,m)+V(m)\bigr).
			\]
			
			Это уравнение утверждает, что минимальная стоимость
			до цели получается выбором лучшего
			узла-преемника и добавлением стоимости перемещения
			к этому узлу.
			
			\medskip
			На каждой итерации мы обновляем оценки для
			всех узлов, используя текущие оценки их
			соседей.
			
			\medskip
			Мы повторяем этот процесс до тех пор, пока оценки
			не сойдутся.
			
		\end{flushleft}
	\end{frame}
	
	%------------------------------------------------
	\begin{frame}{Итерация по значениям: восстановление пути}
		\begin{flushleft}
			
			После сходимости алгоритма мы можем восстановить
			кратчайший путь из любого начального узла $s$.
			
			\medskip
			В каждом узле $k$ мы выбираем преемника $m$,
			минимизирующего сумму стоимости ребра и
			стоимости до цели:
			
			\[
			m=\arg\min_{m\in\mathcal{N}(k)}
			\bigl(c(k,m)+V(m)\bigr).
			\]
			
			Мы переходим к выбранному преемнику и повторяем,
			пока не достигнем целевого узла $g$.
			
		\end{flushleft}
	\end{frame}
	
	
	
	
	\begin{frame}{Литература}
		\begin{itemize}
			
			\item Erik Demaine, Ronald Rivest, Srini Devadas -- MIT:
			\bref{https://ocw.mit.edu/courses/6-006-introduction-to-algorithms-spring-2008/pages/lecture-notes/}
			{Introduction to Algorithms: Lecture Notes, lectures 15--18 (Shortest Paths and Dijkstra).}
			
			\item Jason Ku, Justin Solomon -- MIT:
			\bref{https://ocw.mit.edu/courses/6-006-introduction-to-algorithms-spring-2020/resources/mit6_006s20_lec13/}
			{Introduction to Algorithms: Lecture 13, Dijkstra's Algorithm.}
			
			\item MIT OpenCourseWare:
			\bref{https://ocw.mit.edu/courses/6-231-dynamic-programming-and-stochastic-control-fall-2015/resources/mit6_231f15_lec10/}
			{Dynamic Programming and Stochastic Control: Lecture 10, Bellman's Equation and Value Iteration.}
			
		\end{itemize}
	\end{frame}



\myqrframe




\end{document}
